\documentclass[10pt,a4paper]{amsart}
\usepackage[margin=19mm]{geometry}
\usepackage{amsmath,amssymb,mathtools}
\usepackage[scr=rsfs,cal=euler]{mathalfa}
\usepackage{xfrac}
\usepackage[unicode,hidelinks,bookmarks=false]{hyperref}
\newcommand{\ie}{i.e.\ }
\newcommand{\cf}{cf.\ }
\newcommand{\resp}{respectively\ }
\newcommand{\Subsection}{\subsection}
\providecommand{\nobreakdash}{\nobreak\hbox{-}}
\setlength{\emergencystretch}{2em}
\multlinegap0pt
\usepackage{amssymb}
\usepackage{mathtools}
\usepackage[shortlabels]{enumitem}
\newtheorem{lemma}{Lemma}
\newcommand{\Lalg}{{\mathsf{L}}}
\newcommand{\Mf}{M(\varphi)}          % Verma module of weight varphi
\newcommand{\vf}{v_\varphi}           % highest-weight vector
\title{A Sugawara--Legendre mechanism\\ for the four-point Heisenberg algebra}
\author{Felipe Albino dos Santos}
\date{Source: arXiv:2605.06090v2 (CC BY 4.0)}
\begin{document}
\maketitle

\noindent\emph{Source and licence.} A Sugawara--Legendre mechanism\\ for the four-point Heisenberg algebra. Version arXiv:2605.06090v2 (CC BY 4.0), distributed by arXiv under the Creative Commons Attribution 4.0 International licence, http://creativecommons.org/licenses/by/4.0/, as recorded in the arXiv licence metadata for this exact version and shown on the abstract page https://arxiv.org/abs/2605.06090v2. Full licence notice: \texttt{licenses/arxiv-2605.06090-CC-BY-4.0.txt}.\par\smallskip

\noindent\textit{Editorial scope.} Counted unit: the article's lemma lem:L0-level (source0.tex lines 960--979). The article's complete original proof of it is delivered with it (source0.tex lines 981--1014, one \texttt{proof} environment of 34 lines). No mathematics is written, repaired or re-derived: the statement and the proof are the article's own lines, in the article's own notation, and no retained line is substituted or re-flowed.  Retained uncounted as the source-local context the proof needs: lines 427--431; lines 939--941.  Nothing was omitted from any retained span, so the package carries no omission record.  No retained line contains a citation, so the wrapper carries no bibliography and no bibliography entry.  No retained line contains a figure environment, an image-inclusion command or drawing code, so no image is delivered and the package declares no asset dependency.  Editorial formatting only, in addition: an AMS \texttt{amsart} preamble that restates the article's own declarations and macros used by the retained lines, namely the environments \texttt{lemma} on separate counters, together with the standard packages those lines need.  The retained environments print in this document as Lemma 1 in order of appearance, whereas the article prints the same items with the numbering of its own full document.  The complete native source of the article as distributed in the arXiv e-print for this exact version is bundled unchanged as \texttt{sources/199995/source0.tex}.
\begin{equation}\label{eq:psi-r1}
  [b_k, b_l] \;=\; k\,\delta_{k+l,0}\, \omega_1 \, c,
  \qquad\text{that is}\qquad
  [b_k, b_{-k}] \;=\; k\,\omega_1\,c \quad (k \neq 0),
\end{equation}

\begin{equation}\label{eq:L0-def}
  L_0 \;:=\; \frac{1}{q}\,\Lalg \Big|_{\Mf} .
\end{equation}

\begin{lemma}
\label{lem:L0-level}
Let $v = b_{-n_1}\,b_{-n_2}\,\cdots\,b_{-n_s}\,\vf$ be a PBW monomial
with $n_1\geq n_2\geq\cdots\geq n_s\geq 1$ and write
$n := n_1 + n_2 + \cdots + n_s$. Then, for \emph{every} weight
$\varphi$,
\begin{equation}\label{eq:L0-level}
  \Lalg \cdot v \;=\; q\,n\cdot v .
\end{equation}
If moreover $\varphi$ is Shapovalov-regular, so that $q \neq 0$ and
$L_0 = q^{-1}\Lalg$ is defined \eqref{eq:L0-def}, then
\begin{equation}\label{eq:L0-level-normalised}
  L_0 \cdot v \;=\; n\cdot v .
\end{equation}
Consequently $\Lalg$ acts on $\Mf[-n]$ by the scalar $q\,n$ and $L_0$ by
the scalar $n$, and the spectrum of $L_0$ on $\Mf$ is $\{0,1,2,\dots\}$
with multiplicity $p(n)$ at level $n$, where $p(n)$ is the partition
function. Note that the eigenvalue of $\Lalg$ carries the factor $q$,
which is exactly what the normalisation in \eqref{eq:L0-def} removes.
\end{lemma}

\begin{proof}
We prove \eqref{eq:L0-level}; the normalised form
\eqref{eq:L0-level-normalised} follows by dividing by $q$, which is
legitimate exactly when $\varphi$ is Shapovalov-regular.
We proceed by induction on the length $s$.
For $s=0$ ($v=\vf$), all $b_k\vf = 0$ for $k\geq 1$, so $\Lalg\,\vf = 0
= 0\cdot\vf$.
For $s\geq 1$, write $v = b_{-n_1}\,w$ with $w = b_{-n_2}\cdots b_{-n_s}\vf$
of weight $-(n-n_1)$. Using
$[b_k,b_{-n_1}] = \delta_{k,n_1}\,n_1\,\omega_1\,c$
(the genus-zero coefficients \eqref{eq:psi-r1}, for which
$\psi_{k,-k} = k\,\omega_1$), we compute
\begin{align*}
  L_0 \cdot v
  &= \frac{1}{q}
     \sum_{k\geq 1} b_{-k}\,b_k\,b_{-n_1}\,w \\
  &= \frac{1}{q}
     \sum_{k\geq 1} b_{-k}
       \bigl(b_{-n_1}\,b_k + \delta_{k,n_1}\,n_1\,\omega_1\,c\bigr)\,w \\
  &= b_{-n_1}\,\Bigl(\tfrac{1}{q}\sum_{k\geq 1}b_{-k}b_k\Bigr)\,w
     + \frac{n_1\,q}{q}\,b_{-n_1}\,w \\
  &= b_{-n_1}\,(L_0\,w) \;+\; n_1\,v
   \;=\; (n - n_1)\,v + n_1\,v \;=\; n\,v,
\end{align*}
where we used the inductive hypothesis $L_0\,w = (n-n_1)\,w$ in the last
line, and $c$ acts as the scalar $\varphi(c)$ on $\Mf$ throughout.
This proves \eqref{eq:L0-level}.

The level-$n$ weight space $\Mf[-n]$ is spanned by all such PBW monomials
with $n_1+\cdots+n_s = n$; the count of such monomials is $p(n)$ by the
PBW theorem applied to $\mathcal{H}_2^-$. Each such monomial is a
$L_0$-eigenvector with eigenvalue $n$, so $L_0$ acts as the scalar $n$
on $\Mf[-n]$.
\end{proof}

\end{document}
